\section{Existing guarantees in the units of this paper}
\label{app:comparison}

We separate two tasks: accounting for oracle submissions, and evaluating
the dependence quantities in the representative schedule bounds of
Table~\ref{tab:resource-comparison} and Section~\ref{sec:related-work}.
The latter uses one explicit path, evaluated for TC, DTC, and effective TC.
The comparison substitutes into
published sufficient upper bounds; it is not a lower bound on another
sampler or a tight worst-case complexity over forests.
Appendix~\ref{app:oracle-comparison} records the source-specific oracle
conditions and output objectives.

\subsection{Submissions and parallel depth}
\label{app:comparison-states}

One submitted masked state returns all requested singleton rows; repeated
submissions and different inputs in a parallel batch are charged separately.
A query for one conditional row can therefore be implemented by one such
submission, so a published row-query upper bound is also a full-output
submission upper bound.  This conversion need not be tight.
A product-unmasking schedule with $K$ updates needs at most $K$ submissions
and depth $K$.  Our completed feasible run uses
$\CounterfactualQueries+\CommitRounds+1$ submissions, including preprocessing.
All these counts exclude local arithmetic and vocabulary readout costs.

Parallel verification does not require sequential prefix queries.
Write $\mu$ for the target law used by the external verification sampler,
$h$ for its accepted history, and
$z=(z_1,\ldots,z_b)$ for an ordered candidate block, with
$z_{<j}=(z_1,\ldots,z_{j-1})$.
Once $z$ and $h$ are available, every input in
the chain-rule product
\[
 \mu(z\mid h)=\prod_j\mu(z_j\mid h,z_{<j})
\]
is known and the conditional rows can be queried in parallel.
This is consistent with the sublinear-depth Anari results summarized in
the table; their exact resource statements are recorded only in
Appendix~\ref{app:comparison-exact-oracles}.

\subsection{One path with linear TC and DTC}
\label{app:comparison-path}

All entropies and mutual informations below are under $X\sim P$.
Using the \hyperlink{app-finite-entropy}{finite entropy}
$H(p)=-\sum_xp(x)\log p(x)$, write $H(A):=H(\mathcal L_P(A))$ and,
for finite random variables $A,B,C$,
\[
 \begin{aligned}
 H(A\mid C)&:=\sum_cP(C=c)H(P_{A\mid C=c}),\\
 I(A;B\mid C)&:=\sum_cP(C=c)
   \KLDivergence{P_{AB\mid C=c}}
                {P_{A\mid C=c}\otimes P_{B\mid C=c}}.
 \end{aligned}
\]
The sums run over positive-probability values of $C$;
omitting $C$ gives unconditional mutual information.
Here $P_{AB\mid C=c}$ denotes the conditional joint law and
$P_{A\mid C=c},P_{B\mid C=c}$ its marginals.
With $X_{-i}:=(X_j)_{j\ne i}$, write
\[
 \mathrm{TC}(P):=\sum_iH(X_i)-H(X),\qquad
 \mathrm{DTC}(P):=H(X)-\sum_iH(X_i\mid X_{-i}).
\]
For every forest law, choose a rooted ordering in which each parent
precedes its children.  In the factorization \eqref{app:eq:S0}, a nonroot
vertex $j$, conditional on earlier vertices, has law
$\EndpointKernel{\ParentOf{j}}{j}(\cdot\mid X_{\ParentOf{j}})$:
none of its descendants has appeared yet, and its parent separates it
from the other earlier vertices.  Roots belong to independent components
and retain their marginal laws.  The entropy chain rule therefore gives
\[
 H(X)=\sum_{r\in\TargetRoots}H(X_r)
       +\sum_{j\notin\TargetRoots}H(X_j\mid X_{\ParentOf{j}}).
\]
Subtracting from $\sum_iH(X_i)$ cancels the
root entropies and yields
\begin{align}
 \mathrm{TC}(P)
 &=\sum_{j\notin\TargetRoots}
      [H(X_j)-H(X_j\mid X_{\ParentOf{j}})]\notag\\
 &=\sum_{j\notin\TargetRoots}I(X_j;X_{\ParentOf{j}}).
 \label{app:eq:forest-tc-identity}
\end{align}
A forest need not have large TC: the matching lower-bound witness becomes
weakly coupled as accuracy changes.  The following different witness
has constant dependence per edge while retaining the response and tail
assumptions.  All background tokens share one type, making rare-token
exchanges invisible to the conditional rows.

\paragraph{Intuition.}
The vocabulary has only two dependence-relevant types.  The type chain
has a fixed nonzero dependence between adjacent positions, while token
variation within a type contributes only independent emissions.
Consequently each edge contributes a positive constant to TC, and each
interior position contributes a positive constant to DTC:
\[
 \mathrm{TC}(P)\simeq N,\qquad \mathrm{DTC}(P)\simeq N.
\]
The background vocabulary can grow without changing these dependence
scales.  The proof checks both class membership and the entropy identities.

\begin{proposition}[A shared-class path with linear TC and DTC]
\label{app:prop:path-in-class}
Fix $\RankTailExponent>1$ and put $p=2^{-\RankTailExponent}$.
For every $N,\VocabSize\ge3$, let
$\Vocab=\{a_1,a_2\}\sqcup\mathcal B$, $|\mathcal B|=\VocabSize-2$, with
\[
 \pi(a_1)=1-p,\qquad \pi(a_2)=p/2,\qquad
 \pi(b)=\frac{p}{2(\VocabSize-2)}\quad(b\in\mathcal B).
\]
Define $\tau(a_1)=\mathsf A$ and $\tau(a)=\mathsf B$ for $a\ne a_1$.
The type masses are $P_{\mathsf A}=1-p$, $P_{\mathsf B}=p$.
In the type order $(\mathsf A,\mathsf B)$, set
\begin{equation}
 J=\begin{pmatrix}1-3p/2&p/2\\p/2&p/2\end{pmatrix},
 \qquad
 g(t,t')=\frac{J(t,t')}{P_tP_{t'}},\qquad
 K(y\mid x)=\pi(y)g(\tau(x),\tau(y)).
 \label{app:eq:two-type-kernel}
\end{equation}
Then $P(x)=\pi(x_1)\prod_{j=2}^NK(x_j\mid x_{j-1})$ is a strictly
positive path law with all marginals $\pi$.
There is a constant $\omega_\star(p)>0$, independent of $N,\VocabSize$,
such that this law satisfies (RT\textcolor{red}{; \cref{app:eq:RT}})--(UEN\textcolor{red}{; \cref{app:eq:UEN}}) with
$\ResponseConstant=\ResponseExponent=1$ and every
$\EdgeSignal\le\omega_\star(p)$, and (RF\textcolor{red}{; \cref{app:eq:RF}}) with $\FrequencyConstant=1$.
Its exact oracle satisfies both oracle accuracy conditions.
Moreover,
\begin{equation}
 \mathrm{TC}(P)=\Theta(N),\qquad \mathrm{DTC}(P)=\Theta(N),
 \label{app:eq:path-tc-dtc}
\end{equation}
where the constants depend only on $p$.
Consequently this path belongs to each main-text benchmark class for all
sufficiently large $N$ along that class's vocabulary and accuracy regime.
\end{proposition}

\begin{proof}
\emph{1. Kernel normalization and stationary marginal.}
For a type $t$, the definition of its mass is
$P_t=\sum_{x:\tau(x)=t}\pi(x)$.
Here $0<p<1/2$, and the four entries of $J$ are positive.  Its row sums are
\[
 (1-3p/2)+p/2=1-p=P_{\mathsf A},\qquad
 p/2+p/2=p=P_{\mathsf B}.
\]
Since $J$ is symmetric, its column sums are the same.
For a fixed token $x$ of type $t=\tau(x)$, group the sum over $y$ by type:
\begin{align*}
 \sum_yK(y\mid x)
 &=\sum_{t'\in\{\mathsf A,\mathsf B\}}
       \sum_{y:\tau(y)=t'}\pi(y)\frac{J(t,t')}{P_tP_{t'}}\\
 &=\sum_{t'\in\{\mathsf A,\mathsf B\}}
       P_{t'}\frac{J(t,t')}{P_tP_{t'}}
 =\frac1{P_t}\sum_{t'}J(t,t')=1.
\end{align*}
Thus $K$ is a probability kernel.  For a fixed $y$ of type $t'$,
the column-sum identity similarly yields
\begin{align*}
 \sum_x\pi(x)K(y\mid x)
 &=\pi(y)\sum_{t\in\{\mathsf A,\mathsf B\}}
       \sum_{x:\tau(x)=t}\pi(x)\frac{J(t,t')}{P_tP_{t'}}\\
 &=\pi(y)\frac1{P_{t'}}\sum_tJ(t,t')=\pi(y).
\end{align*}
Starting the path from $\pi$ therefore gives marginal $\pi$ at every
position, by induction along the path.  Also,
\[
 \pi(x)K(y\mid x)
 =\pi(x)\pi(y)\frac{J(\tau(x),\tau(y))}
                      {P_{\tau(x)}P_{\tau(y)}}
 =\pi(y)K(x\mid y),
\]
which verifies endpoint compatibility.  Positivity of $\pi$ and $K$
gives strict positivity of the path law.

\emph{2. Boundary-conditioned rows and the edge signal.}
Substituting the kernel formula into the path law gives
\[
 P(x)=\left(\prod_{j=1}^N\pi(x_j)\right)
      \left(\prod_{j=2}^Ng(\tau(x_{j-1}),\tau(x_j))\right).
\]
For fixed values outside readout $j$, all factors not incident to $j$
cancel.  Write $\mathcal N_F(j)=\{j-1,j+1\}\cap[N]$ for its path
neighbors.  Its normalized conditional row is
\[
 P(X_j=y\mid X_{-j}=x_{-j})
 =\frac{\pi(y)\prod_{k\in\mathcal N_F(j)}
                   g(\tau(y),\tau(x_k))}
        {\sum_z\pi(z)\prod_{k\in\mathcal N_F(j)}
                   g(\tau(z),\tau(x_k))}.
\]
Within type $t$, every factor $g$ is constant, so the conditional token
law given the type is always $\pi(y)/P_t$ for $\tau(y)=t$.
Summing the numerator separately over the two types shows that the
conditional odds of type $\mathsf A$ against type $\mathsf B$ are
\begin{align*}
 &\frac{P(\tau(X_j)=\mathsf A\mid X_{-j}=x_{-j})}
        {P(\tau(X_j)=\mathsf B\mid X_{-j}=x_{-j})}\\
 &\qquad=
 \frac{P_{\mathsf A}\prod_{k\in\mathcal N_F(j)}g(\mathsf A,\tau(x_k))}
      {P_{\mathsf B}\prod_{k\in\mathcal N_F(j)}g(\mathsf B,\tau(x_k))}
 =\rho\prod_{k\in\mathcal N_F(j)}r_{\tau(x_k)},\\
 &\rho:=\frac{P_{\mathsf A}}{P_{\mathsf B}}=\frac{1-p}{p},
 \qquad r_t:=\frac{g(\mathsf A,t)}{g(\mathsf B,t)}
 \quad(t\in\{\mathsf A,\mathsf B\}).
\end{align*}
Thus $\rho$ is the marginal type-odds ratio, and $r_t$ is the
multiplicative contribution of a neighbor of type $t$.
Their values follow from the four entries
\[
 g(\mathsf A,\mathsf A)=\frac{1-3p/2}{(1-p)^2},\quad
 g(\mathsf A,\mathsf B)=g(\mathsf B,\mathsf A)=\frac1{2(1-p)},\quad
 g(\mathsf B,\mathsf B)=\frac1{2p}:
\]
\begin{align*}
 \PathOddsA&=
 \frac{g(\mathsf A,\mathsf A)}{g(\mathsf B,\mathsf A)}
 =
 \frac{(1-3p/2)/(1-p)^2}{1/[2(1-p)]}
 =\frac{2-3p}{1-p},\\
 \PathOddsB&=
 \frac{g(\mathsf A,\mathsf B)}{g(\mathsf B,\mathsf B)}
 =
 \frac{1/[2(1-p)]}{1/(2p)}
 =\frac{p}{1-p}.
\end{align*}
In particular, $\PathOddsA-\PathOddsB=2(1-2p)/(1-p)>0$.

Fix a directed edge $i\to j$ and a complete boundary
$z\in\Vocab^{[N]\setminus\{i,j\}}$.
The two rows to compare are precisely
\begin{align*}
 \BoundaryRow{i}{j}{a_1}{z}
 &=\mathcal L_P(X_j\mid X_i=a_1,X_{[N]\setminus\{i,j\}}=z),\\
 \BoundaryRow{i}{j}{a_2}{z}
 &=\mathcal L_P(X_j\mid X_i=a_2,X_{[N]\setminus\{i,j\}}=z).
\end{align*}
Here $\tau(a_1)=\mathsf A$ and $\tau(a_2)=\mathsf B$.
Put $f(o)=o/(1+o)$, the type-$\mathsf A$ probability at odds $o$, and
define the contribution of the other neighbors by
\[
 q:=\prod_{k\in\mathcal N_F(j)\setminus\{i\}}r_{\tau(z_k)}
 \in\{1,\PathOddsA,\PathOddsB\}.
\]
The empty product is one; a path readout has at most one other neighbor.
The type-$\mathsf A$ weights of these two rows are, respectively,
\begin{align*}
 u&:=P(\tau(X_j)=\mathsf A\mid X_i=a_1,X_{[N]\setminus\{i,j\}}=z)
       =f(\rho\PathOddsA q),\\
 u'&:=P(\tau(X_j)=\mathsf A\mid X_i=a_2,X_{[N]\setminus\{i,j\}}=z)
       =f(\rho\PathOddsB q).
\end{align*}
Within each type, both rows have the same token law $\pi(y)/P_t$.
Expanding TV over the disjoint type supports therefore gives
\begin{align*}
 &\dTV\bigl(\BoundaryRow{i}{j}{a_1}{z},
             \BoundaryRow{i}{j}{a_2}{z}\bigr)\\
 &\quad=\frac12\left(
    \sum_{\tau(y)=\mathsf A}|u-u'|\frac{\pi(y)}{P_{\mathsf A}}
   +\sum_{\tau(y)=\mathsf B}|(1-u)-(1-u')|
                            \frac{\pi(y)}{P_{\mathsf B}}\right)\\
 &\quad=\frac12\bigl(|u-u'|+|(1-u)-(1-u')|\bigr)
 =u-u'=:D(q),\\
 &D(q)=f(\rho\PathOddsA q)-f(\rho\PathOddsB q)
 =\frac{\rho q(\PathOddsA-\PathOddsB)}
        {(1+\rho\PathOddsA q)(1+\rho\PathOddsB q)}>0.
\end{align*}
The two sums of $\pi(y)/P_t$ equal one by the definition of $P_t$.
The conditional row depends on the source token only through its type,
so $a_1,a_2$ attain the full directed response:
\[
 \DirectedResponse{i}{j}{z}
 =\max_{c,c'\in\Vocab}
   \dTV\bigl(\BoundaryRow{i}{j}{c}{z},
              \BoundaryRow{i}{j}{c'}{z}\bigr)
 =D(q).
\]
Define
$\omega_\star(p):=\min\{D(1),D(\PathOddsA),D(\PathOddsB)\}>0$.
This minimum depends only on $p$, not on $N,\VocabSize$ or the boundary.
The same calculation applies to either orientation of every path edge,
so the required condition \eqref{app:eq:UEN} is
\[
 \min\bigl\{\DirectedResponse{i}{j}{z},
             \DirectedResponse{j}{i}{z}\bigr\}
 \ge\omega_\star(p)\ge\EdgeSignal
 \qquad\text{(UEN\textcolor{red}{; \cref{app:eq:UEN}})}.
\]

\emph{3. Tail response and the rank--frequency envelope.}
Recall that $\PosMarginal{i}=\pi$ at every position, so the marginal
tail in \eqref{app:eq:tail} is
$\TailSet{i}{t}=\{c\in\Vocab:\pi(c)\le t\}$.
For a nonedge $\{i,j\}\notin\TargetEdges$, the conditional row does
not involve the source value, and hence for every $c,c'\in\Vocab$,
\[
 \BoundaryRow{i}{j}{c}{z}=\BoundaryRow{i}{j}{c'}{z},
 \qquad
 \dTV\bigl(\BoundaryRow{i}{j}{c}{z},
            \BoundaryRow{i}{j}{c'}{z}\bigr)=0.
\]
For an edge, the two threshold ranges are as follows.
\begin{itemize}
 \item If $t<1-p=\pi(a_1)$, then
 $\TailSet{i}{t}\subseteq\{a_2\}\cup\mathcal B=\tau^{-1}(\mathsf B)$.
 Thus all permitted source tokens have the same type and give the same
 readout row; their maximum TV distance is zero.
 \item If $t\ge1-p$, then $\TailSet{i}{t}=\Vocab$, since $1-p$ is the
 largest marginal mass.  The full directed response is $D(q)$ from
 step~2, so it remains to prove $D(q)\le1-p$.
\end{itemize}

Here is the needed bound with its maximization made explicit.
Set $x=\sqrt{\PathOddsA/\PathOddsB}>1$ and
$\zeta=\rho\PathOddsB q>0$.  Then
\[
 D(q)=\frac{x^2\zeta}{1+x^2\zeta}-\frac\zeta{1+\zeta}
     =\frac{(x^2-1)\zeta}{(1+x^2\zeta)(1+\zeta)}.
\]
The denominator is $1+(1+x^2)\zeta+x^2\zeta^2$.
The quotient rule gives
\[
 \frac{d}{d\zeta}
 \frac{(x^2-1)\zeta}{(1+x^2\zeta)(1+\zeta)}
 =\frac{(x^2-1)(1-x^2\zeta^2)}
        {(1+x^2\zeta)^2(1+\zeta)^2}.
\]
The derivative is positive for $\zeta<1/x$ and negative for $\zeta>1/x$.
Substituting $\zeta=1/x$ therefore gives
\[
 \sup_{q>0}D(q)
 =\frac{(x^2-1)/x}{(1+x)(1+1/x)}
 =\frac{x-1}{x+1}.
\]
All quantities are positive, and
\[
 \frac{x-1}{x+1}\le1-p
 \quad\Longleftrightarrow\quad
 px\le2-p
 \quad\Longleftrightarrow\quad
 p^2x^2\le(2-p)^2.
\]
Finally,
\[
 p^2x^2=p(2-3p),\qquad
 (2-p)^2-p(2-3p)=4-6p+4p^2>0
 \quad(0<p<1/2).
\]
Combining the nonedge and edge cases gives, for every distinct $i,j$,
complete boundary $z$, and $t\in[0,1]$,
\[
 \begin{aligned}
 &\max_{c,c'\in\TailSet{i}{t}}
 \dTV\bigl(\BoundaryRow{i}{j}{c}{z},
            \BoundaryRow{i}{j}{c'}{z}\bigr)\\
 &\qquad\le
 \begin{cases}
  0,&0\le t<1-p,\\
  1-p,&1-p\le t\le1
 \end{cases}
 \le t=\ResponseConstant t^{\ResponseExponent}
 \qquad\text{(RT\textcolor{red}{; \cref{app:eq:RT}})},
 \end{aligned}
\]
where $\ResponseConstant=\ResponseExponent=1$, and the maximum is zero
for a tail with at most one token, as in \eqref{app:eq:RT}.

For (RF\textcolor{red}{; \cref{app:eq:RF}}), $a_1$ has mass $1-p>1/2$ and is the largest atom.
The mass of $a_2$ is $p/2=2^{-\RankTailExponent-1}\le2^{-\RankTailExponent}$.
Each background mass obeys, since $\VocabSize\ge3$,
\[
 \frac{p}{2(\VocabSize-2)}
 \le\frac1{4(\VocabSize-2)}
 \le\frac1{\VocabSize};
 \qquad
 \VocabSize\le4(\VocabSize-2)
 \ \Longleftrightarrow\ 3\VocabSize\ge8.
\]
The background masses are no larger than $p/2$.
Therefore, with the ranked tokens of \eqref{app:eq:hidden-ranks}, every
position $i$ and rank $k\in[\VocabSize]$ satisfy
\begin{align*}
 \PosMarginal{i}(\RankedToken{i}{k})
 &\le
 \begin{cases}
  1,&k=1,\\
  2^{-\RankTailExponent},&k=2,\\
  \VocabSize^{-1},&3\le k\le\VocabSize
 \end{cases}
 \\
 &\le k^{-\RankTailExponent}+\VocabSize^{-1}
 =\FrequencyConstant(k^{-\RankTailExponent}+\VocabSize^{-1})
 \qquad\text{(RF\textcolor{red}{; \cref{app:eq:RF}})},
\end{align*}
with $\FrequencyConstant=1$.
If $\VocabSize=3$, the possible tie at rank two has mass $p/2$ on
either token, so the same bounds apply under the public tie-breaking order.
This is the finite envelope \eqref{app:eq:RF}; the asymptotic part of
(RF\textcolor{red}{; \cref{app:ass:rate-regime}}) is supplied by the prescribed sequence
$\VocabSize=N^{\VocabGrowthExponent+o(1)}$.
The exact oracle has zero Hellinger and KL error.

\emph{4. Total correlation.}
Let $T_j=\tau(X_j)$.  Its transition probability follows by grouping $K$:
\[
 P(T_{j+1}=t'\mid T_j=t)
 =\sum_{\tau(y)=t'}\pi(y)g(t,t')
 =\frac{J(t,t')}{P_t}.
\]
Since its marginal is $P_t$, the adjacent type joint law is $J$.
The adjacent token joint law is
\[
 P(X_j=x,X_{j+1}=y)
 =J(t,t')\frac{\pi(x)}{P_t}\frac{\pi(y)}{P_{t'}},
 \qquad t=\tau(x),\quad t'=\tau(y).
\]
Its likelihood ratio relative to $\pi(x)\pi(y)$ is
$J(t,t')/(P_tP_{t'})$, which depends only on the types.
Grouping the mutual-information sum by type therefore gives
\[
 I(X_j;X_{j+1})
 =\sum_{t,t'}J(t,t')\log\frac{J(t,t')}{P_tP_{t'}}
 =I(T_j;T_{j+1})=:c.
\]
\hypertarget{app-positive-pair-mi}{}
To see why this constant is positive, the displayed sum is
$\KLDivergence{J}{(P_tP_{t'})_{t,t'}}$.
KL is zero exactly when its two laws agree.  Such equality would make
$J$ a product table and hence rank one, but
\[
 \det J=(1-3p/2)(p/2)-(p/2)^2
       =\frac{p(1-2p)}2>0.
\]
It depends only on $p$, so \eqref{app:eq:forest-tc-identity} gives
$\mathrm{TC}(P)=(N-1)c=\Theta(N)$.

\emph{5. Dual total correlation.}
The chain rule and conditional Markov property give, respectively,
\[
 H(X)=H(X_1)+\sum_{i=2}^NH(X_i\mid X_{i-1})
\]
and
\[
 \sum_iH(X_i\mid X_{-i})
 =H(X_1\mid X_2)
  +\sum_{i=2}^{N-1}H(X_i\mid X_{i-1},X_{i+1})
  +H(X_N\mid X_{N-1}).
\]
Subtracting cancels the final-coordinate term and yields
\begin{align}
 \mathrm{DTC}(P)
 &=H(X_1)-H(X_1\mid X_2)\notag\\
 &\quad+\sum_{i=2}^{N-1}
   \bigl[H(X_i\mid X_{i-1})
         -H(X_i\mid X_{i-1},X_{i+1})\bigr]\notag\\
 &=I(X_1;X_2)+\sum_{i=2}^{N-1}
                   I(X_i;X_{i+1}\mid X_{i-1}).
 \label{app:eq:path-dtc-identity}
\end{align}
To evaluate an interior term, fix the previous type $t$.
The conditional joint table of $(T_i,T_{i+1})$ is
\[
 P(T_i=u,T_{i+1}=v\mid T_{i-1}=t)
 =\frac{J(t,u)}{P_t}\frac{J(u,v)}{P_u}.
\]
This is obtained from $J$ by multiplying row $u$ by the positive factor
$J(t,u)/(P_tP_u)$.  Its determinant is consequently
\[
 \frac{J(t,\mathsf A)J(t,\mathsf B)}
      {P_t^2P_{\mathsf A}P_{\mathsf B}}\det J>0.
\]
The two conditional type variables are therefore dependent for either
$t$, and their conditional mutual information is a positive constant
depending only on $p$.

It remains to relate this type calculation to the token term.
Conditioning on $X_{i-1}=x$ affects the future law only through $\tau(x)=t$.
For this calculation write $\Pr_t(\cdot)=\Pr(\cdot\mid T_{i-1}=t)$.
For tokens $y,z$ of types $u,v$, the emission factors cancel as follows:
\begin{align*}
 &\frac{\Pr_t(X_i=y,X_{i+1}=z)}
        {\Pr_t(X_i=y)\Pr_t(X_{i+1}=z)}\\
 &\quad=
 \frac{\Pr_t(T_i=u,T_{i+1}=v)\,[\pi(y)/P_u]\,[\pi(z)/P_v]}
      {\Pr_t(T_i=u)\,[\pi(y)/P_u]\,
       \Pr_t(T_{i+1}=v)\,[\pi(z)/P_v]}\\
 &\quad=
 \frac{\Pr_t(T_i=u,T_{i+1}=v)}
      {\Pr_t(T_i=u)\Pr_t(T_{i+1}=v)}.
\end{align*}
Grouping the token sum by $(u,v)$ therefore gives
\begin{align*}
 &I(X_i;X_{i+1}\mid T_{i-1}=t)\\
 &=\sum_{u,v}\Pr_t(T_i=u,T_{i+1}=v)
   \log\frac{\Pr_t(T_i=u,T_{i+1}=v)}
             {\Pr_t(T_i=u)\Pr_t(T_{i+1}=v)}\\
 &\qquad\times
   \underbrace{\left(\sum_{y:\tau(y)=u}\frac{\pi(y)}{P_u}\right)}_{=1}
   \underbrace{\left(\sum_{z:\tau(z)=v}\frac{\pi(z)}{P_v}\right)}_{=1}\\
 &=I(T_i;T_{i+1}\mid T_{i-1}=t).
\end{align*}
The two sums equal one because $P_u=\sum_{\tau(y)=u}\pi(y)$ and similarly
for $P_v$.  The conditional future law is the same for every previous token
of type $t$, so averaging over the previous token, or equivalently over its
type with weights $P_t$, yields
\[
 I(X_i;X_{i+1}\mid X_{i-1})
 =I(T_i;T_{i+1}\mid T_{i-1}).
\]
Every interior term is the same positive constant by stationarity and is
at most $\log2$, since the types are binary.
Together with $I(X_1;X_2)=c>0$, the identity above proves
$\mathrm{DTC}(P)=\Theta(N)$.

All constants used for the signal and dependence depend only on
$p=2^{-\RankTailExponent}$.  Each main benchmark has
$\EdgeSignal\to0$, so eventually $\EdgeSignal\le\omega_\star(p)$;
its prescribed vocabulary sequence supplies the remaining asymptotic
part of (RF\textcolor{red}{; \cref{app:ass:rate-regime}}).  This proves membership in those benchmark classes.
\end{proof}

\paragraph{What the substitution establishes.}
For this path, the TC/DTC sufficient counts listed in
Table~\ref{tab:resource-comparison} become
$\widetilde O(1+N/\varepsilon)$ at forward-KL tolerance $\varepsilon$.
These displayed guarantees alone do not yield a sublinear count at fixed
accuracy.  They do not establish that those samplers require linear work:
exact singleton updates already give an $N$-call cap.
Our case-(ii) benchmark also applies, with the balanced count in
\eqref{main:eq:balanced-upper} on its sublinear range.
The calculation concerns one class member, not every forest or every
dependence-adaptive schedule.

\begin{proposalcontext}{Other dependence measures}
No corresponding quantitative comparison is made here for entropy-based
bounds or refined unmasking growth complexity.
\end{proposalcontext}

\subsection{Effective total correlation on the same path}
\label{app:comparison-effective-tc}

To compare with \citet{dmitriev2026efficient}, let $Y(t)$ be an
independently masked copy of $X\sim P$: each $Y_i(t)$ equals $X_i$
with probability $e^{-t}$ and equals $\MASK$ otherwise.
Write $Y_{-(i,j)}(t)$ for all coordinates except $i,j$.
Their effective total correlation, written $\EffectiveTC$ here, is
\begin{equation}
 \EffectiveTC(P)
 :=\int_0^\infty\min\{1,t\}\MaskedPairInformation_P(t)\,dt,
 \qquad
 \MaskedPairInformation_P(t)
 :=\sum_{i\ne j}I\bigl(Y_i(t);Y_j(t)\mid Y_{-(i,j)}(t)\bigr).
 \label{app:eq:effective-tc-definition}
\end{equation}
These are equation~(16) of
\href{https://arxiv.org/html/2602.15008v2}{arXiv:2602.15008v2}.
Its Lemma~16 gives
$\EffectiveTC(P)\le\min\{\mathrm{TC}(P),\mathrm{DTC}(P)\}$.
Large TC and DTC alone do not imply large effective TC: that paper's
Proposition~5 supplies a counterexample. We therefore evaluate
$\EffectiveTC$ on our path directly.

\paragraph{Intuition.}
For an interior edge, reveal its two endpoints and their two outer
neighbors. This occurs with probability $e^{-4t}$.
Fixing those neighbors isolates a dependent two-vertex conditional law,
regardless of other reveals. Summing over the $N-3$ interior edges gives
$\MaskedPairInformation_P(t)\gtrsim_{\RankTailExponent}Ne^{-4t}$;
integrating over a fixed time interval keeps a linear contribution.

\begin{corollary}[Linear effective TC on the comparison path]
\label{app:cor:path-effective-tc}
For the path law in Proposition~\ref{app:prop:path-in-class},
with fixed $\RankTailExponent>1$,
$\EffectiveTC(P)=\Theta(N)$ as $N\to\infty$.
The constants depend only on $\RankTailExponent$, not on $\VocabSize$.
\end{corollary}

\begin{proof}
% Please keep the mask definitions and their chain-rule expansion on one page.
\begin{samepage}
\emph{1. Expand over the visible coordinates.}
Take $N\ge4$, fix $t>0$, and put $\lambda=e^{-t}$.
For a pair $i\ne j$, the visible set outside the pair is
$S\subseteq[N]\setminus\{i,j\}$ with probability
$\lambda^{|S|}(1-\lambda)^{N-2-|S|}$.
Put $M_k:=\mathbf 1\{Y_k(t)\ne\MASK\}$, so the $M_k$ are independent
$\operatorname{Ber}(\lambda)$ variables, independent of $X$.
Fix an exterior mask pattern with visible set $S$ and revealed values
$X_S=x_S$.  The conditioning event is
\[
 \mathcal C=\{M_k=\mathbf 1\{k\in S\}\ \text{for every }k\notin\{i,j\},
                \ X_S=x_S\}.
\]
For this calculation abbreviate $Y_k(t)$ to $Y_k$.
Since $M_i$ is determined by $Y_i$ and $M_j$ by $Y_j$, two applications
of the conditional chain rule give
\begin{align*}
 I(Y_i;Y_j\mid\mathcal C)
 &=I(M_i;Y_j\mid\mathcal C)
   +I(Y_i;Y_j\mid M_i,\mathcal C)\\
 &=I(Y_i;M_j\mid M_i,\mathcal C)
   +I(Y_i;Y_j\mid M_i,M_j,\mathcal C)\\
 &=\lambda^2 I(X_i;X_j\mid X_S=x_S).
\end{align*}
Here $I(M_i;Y_j\mid\mathcal C)=0$ and
$I(Y_i;M_j\mid M_i,\mathcal C)=0$ by mask independence.
The last conditional mutual information averages over $(M_i,M_j)$.
It is zero whenever either endpoint is masked, because that $Y$ value is
constant; when both are visible it equals
$I(X_i;X_j\mid X_S=x_S)$, with weight $\lambda^2$.
Averaging over $x_S$ and then over exterior mask patterns gives
\begin{equation}
 I\bigl(Y_i(t);Y_j(t)\mid Y_{-(i,j)}(t)\bigr)
 =\lambda^2
   \sum_{S\subseteq[N]\setminus\{i,j\}}
    \lambda^{|S|}(1-\lambda)^{N-2-|S|}
    I(X_i;X_j\mid X_S).
 \label{app:eq:masked-pair-information}
\end{equation}
\end{samepage}

\emph{2. Retain a positive dependence across each interior edge.}
Recall the two types $T_i=\tau(X_i)$ and their transition matrix
$\bar K(a,b)=J(a,b)/P_a$ from the proof of
Proposition~\ref{app:prop:path-in-class}.
All entries are positive, and
$\det\bar K=\det J/(P_{\mathsf A}P_{\mathsf B})>0$.
For an interior pair $i,i+1$, where $2\le i\le N-2$, fix the outer
types $T_{i-1}=a,T_{i+2}=b$.
The conditional two-by-two table is
\begin{align*}
 &\Pr(T_i=u,T_{i+1}=v\mid T_{i-1}=a,T_{i+2}=b)
 =\frac{\bar K(a,u)\bar K(u,v)\bar K(v,b)}{Z_{a,b}},\\
 &Z_{a,b}:=\sum_{u,v}\bar K(a,u)\bar K(u,v)\bar K(v,b)>0.
\end{align*}
Here and below type indices range over $\{\mathsf A,\mathsf B\}$.
This table is obtained from $\bar K$ by multiplying row $u$ by
$\bar K(a,u)$, column $v$ by $\bar K(v,b)$, and the whole table
by $Z_{a,b}^{-1}$. Its determinant is
\[
 \frac{\bar K(a,\mathsf A)\bar K(a,\mathsf B)
       \bar K(\mathsf A,b)\bar K(\mathsf B,b)}
      {Z_{a,b}^{\,2}}\det\bar K>0.
\]
Thus the conditional pair has positive mutual information for each
boundary type pair, by the \hyperlink{app-positive-pair-mi}{KL and rank-one
criterion} used above.  There are only four such pairs.  By stationarity,
the positive constant
\[
 c_{\mathrm{pair}}
 :=\min_{a,b}I(T_2;T_3\mid T_1=a,T_4=b)>0
\]
depends only on $p=2^{-\RankTailExponent}$.

If $\{i-1,i+2\}\subseteq S\subseteq[N]\setminus\{i,i+1\}$,
path factorization shows that other
observed coordinates affect only factors outside this pair.
Conditional on $X_S=x_S$, the type pair therefore has the displayed
table with $a=\tau(x_{i-1})$, $b=\tau(x_{i+2})$.
Applying data processing to the two type maps at each $x_S$ and then
averaging gives
\begin{align*}
 I(X_i;X_{i+1}\mid X_S)
 &=\sum_{x_S}P(X_S=x_S)
       I(X_i;X_{i+1}\mid X_S=x_S)\\
 &\ge\sum_{x_S}P(X_S=x_S)
       I(T_i;T_{i+1}\mid X_S=x_S)\\
 &\ge\sum_{x_S}P(X_S=x_S)c_{\mathrm{pair}}
 =c_{\mathrm{pair}}.
\end{align*}
The total mask weight of these sets is
\[
 \sum_{\substack{S\subseteq[N]\setminus\{i,i+1\}\\
                 \{i-1,i+2\}\subseteq S}}
   \lambda^{|S|}(1-\lambda)^{N-2-|S|}
 =\lambda^2\bigl(\lambda+(1-\lambda)\bigr)^{N-4}
 =\lambda^2.
\]
The first $\lambda^2$ requires the two outer neighbors to be visible;
each remaining coordinate contributes $\lambda+(1-\lambda)=1$.
Substituting into \eqref{app:eq:masked-pair-information} yields
$I(Y_i(t);Y_{i+1}(t)\mid Y_{-(i,i+1)}(t))
 \ge c_{\mathrm{pair}}\lambda^4$.

\emph{3. Sum the edges and integrate.}
There are $N-3$ such edges and two orientations per edge in
$\MaskedPairInformation_P(t)$. Nonnegativity of the other terms gives
\[
 \MaskedPairInformation_P(t)\ge
 2(N-3)c_{\mathrm{pair}}e^{-4t}.
\]
Hence
\[
 \EffectiveTC(P)\ge
 2(N-3)c_{\mathrm{pair}}
 \int_0^\infty\min\{1,t\}e^{-4t}\,dt
 =\Omega_{\RankTailExponent}(N).
\]
The integral is a positive finite absolute constant.
The reverse bound follows from
$\EffectiveTC(P)\le\mathrm{DTC}(P)=O_{\RankTailExponent}(N)$,
using the cited Lemma~16 and Proposition~\ref{app:prop:path-in-class}.
\end{proof}
